\section{Proofs for sparse regression}
\label{app:regression}

\subsection{Gaussian estimates and planted residuals}
\label{app:regression:tools}

We use the following standard concentration inequalities. For
$Z\sim N(0,1)$, $t>0$, and the standard normal density $\phi$,
\begin{equation}\label{regression:gaussian-tail}
 \Pbb\{|Z|>t\}\leq e^{-t^2/2},\qquad
 \frac{t}{1+t^2}\phi(t)\leq\overline\Phi(t)\leq\frac{\phi(t)}t.
\end{equation}
\citet[Lemma~1 and (4.3)--(4.4), p.~1325]{laurentMassart2000ChiSquare}
give, for $Q_d\sim\chi_d^2$ and $t>0$,
\begin{equation}\label{regression:chi-tail}
 \Pbb\{Q_d>d+2\sqrt{dt}+2t\}\leq e^{-t},\qquad
 \Pbb\{Q_d<d-2\sqrt{dt}\}\leq e^{-t}.
\end{equation}
For a $d$ by $m$ standard Gaussian matrix $A$, $d\geq m$, each of
\begin{equation}\label{regression:singular-tail}
 s_{\max}(A)\leq\sqrt d+\sqrt m+t,\qquad
 s_{\min}(A)\geq\sqrt d-\sqrt m-t
\end{equation}
fails with probability at most $e^{-t^2/2}$.
These bounds follow from
\citet[Theorem~II.13, p.~353]{davidsonSzarek2001LocalOperatorTheory}
with the correction in
\citet[item~5, p.~1819]{davidsonSzarek2003Corrigendum}.
Their matrix has entry variance $1/d$; multiplying it by $\sqrt d$
gives the standard Gaussian normalization used here.
We square the latter lower bound only when its right side is
nonnegative. Finally, if $N$ is binomial with mean $\mu$, then
$\Pbb\{N>2\mu+3t\}\leq e^{-t}$ and
$\Pbb\{N<\mu/2\}\leq e^{-\mu/8}$.

One additional estimate will be useful. For independent standard
Gaussian variables $Z_j$ and $q>0$, put
$T_j=(|Z_j|-q)_+^2$. Completing the square gives
\begin{align}
 m(q)&=\E e^{T_j/4}-1
       =2\sqrt2e^{q^2/2}\overline\Phi(\sqrt2q)-2\overline\Phi(q)
       \leq 2\phi(q)/q,\\
 \Pbb\left\{\sum_{j=1}^NT_j>4[Nm(q)+t]\right\}&\leq e^{-t}.
 \label{regression:truncated-tail}
\end{align}
Indeed, on $z>q$,
$-z^2/2+(z-q)^2/4=-(z+q)^2/4+q^2/2$.
Integrating over the two tails gives the first identity and
\eqref{regression:gaussian-tail} gives its bound. Markov's inequality,
$\E e^{\sum_jT_j/4}=(1+m(q))^N\leq e^{Nm(q)}$, gives the second line.

\begin{lemma}[Uniform planted fits]\label{regression:fit-concentration}
Under \cref{regression:easy-assumption}, with $S=S^*$,
\begin{equation}\label{regression:fit-concentration-eq}
 \max_{i\in S}\left|\frac{|\widehat\beta_i|}{b_\lambda}-1\right|
 =o_{\Pbb}(1),\qquad
 \frac{\|r_S\|^2}{\omega_\lambda^2}=1+o_{\Pbb}(1).
\end{equation}
These errors can be bounded by deterministic quantities
$O(\sqrt{k/n}+\sqrt{\log p/n}+\log^{-2}p)$ on events whose
probabilities tend to one.
\end{lemma}
\begin{proof}
Put $G=X_S^\top X_S$, $\theta=\sqrt{k/n}+2\sqrt{\log p/n}$,
and $v=M_S^{-1}X_S\beta_S^*$. By
\eqref{regression:singular-tail}, with probability tending to one,
the eigenvalues of $G$ are $n(1+O(\theta))$.
The push-through identity gives
\[
 v=\lambda X_S(G+\lambda I)^{-1}\beta_S^*,\qquad
 \|v\|^2=\lambda^2(\beta_S^*)^\top
              (G+\lambda I)^{-1}G(G+\lambda I)^{-1}\beta_S^*.
\]
Hence
$\|v\|^2=(1+O(\theta))kb^2\lambda^2n/(n+\lambda)^2$.
On the orthogonal complement of $\operatorname{range}X_S$,
$M_S^{-1}$ is the identity; on the column space its eigenvalues
are between zero and one. Conditional chi-square concentration gives
\[
 n-k-O(\sqrt{n\log p})\leq\|M_S^{-1}w\|^2
          \leq n+O(\sqrt{n\log p}+\log p).
\]
Conditional on $X_S$, $v^\top M_S^{-1}w$ is centered normal with
standard deviation at most $\|v\|$. Its magnitude is at most
$\|v\|\sqrt{2\log p}$ with probability tending to one. Since
$r_S=v+\sigma M_S^{-1}w$, and
$2\sigma\|v\|\sqrt{2\log p}\leq
\sqrt{2\log p/n}(n\sigma^2+\|v\|^2)$, this proves the residual estimate.

For the coefficient estimate, remove coordinate $i$ and write
$A_i=M_{S\setminus\{i\}}^{-1}$ and
$y_{(i)}=y-\beta_i^*x_i$. Eliminating the remaining coefficients gives
\begin{equation}\label{regression:loo}
 \widehat\beta_i=\frac{\beta_i^*q_i+\xi_i}{\lambda+q_i},
 \qquad q_i=x_i^\top A_ix_i,\quad \xi_i=x_i^\top A_iy_{(i)}.
\end{equation}
The matrix $A_i$ lies between the orthogonal complement projector
and the identity. By \eqref{regression:chi-tail} and a union bound,
simultaneously in $i\in S$,
$q_i=n+O(k+\sqrt{n\log p}+\log p)$.
Conditional on $(X_{S\setminus\{i\}},w)$, $\xi_i$ is normal with
variance $\|A_iy_{(i)}\|^2$. The preceding signal calculation,
applied to a subset, and eigenvalue interlacing bound
\[
 \|A_iy_{(i)}\|\leq
 Cb\sqrt{k}\lambda/\sqrt n+\sigma\|w\|.
\]
Gaussian tails and a union bound then show that the noise term in
\eqref{regression:loo}, relative to $b_\lambda$, is at most
\[
 C\left[\frac{\sigma}{b}\sqrt{\frac{\log p}{n}}+
              \frac{\lambda}{n}\sqrt{\frac{k\log p}{n}}\right].
\]
The change from $bq_i/(\lambda+q_i)$ to $b_\lambda$ is bounded
relatively by $C\theta\lambda/n$. Under
\eqref{regression:easy-regime}, all these errors tend to zero and
obey the stated bound. The union bounds have failure $O((k+1)/p)$,
which tends to zero because $k/p=O(1/\log p)$.
\end{proof}

\subsection{Proof of the perspective thresholds}
\label{app:regression:easy}

\begin{proof}[Proof of \cref{regression:easy-thresholds}(a)]
Condition on $(X_S,w)$ and put $a=X^\top r_S$,
$m_0=\lambda\min_{i\in S}|\widehat\beta_i|$.
The null variables $Z_j=a_j/\|r_S\|$, $j\notin S$, are independent
standard normal variables. By \cref{regression:fit-concentration},
$m_0/\|r_S\|=(1+o_{\Pbb}(1))\tau_\lambda$.
For the positive assertion, a union bound gives
\[
 \Pbb\{\max_{j\notin S}|Z_j|>
                  \sqrt{2\log p+2\log\log p}\}\leq1/\log p.
\]
The selected threshold is larger with probability tending to one;
\cref{regression:exact-condition} proves all conclusions.
For the negative assertion, take
$t=\sqrt{(2-\delta/2)\log p}$. On the high-probability fit event,
$t>m_0/\|r_S\|$. Moreover,
\[
 \Pbb\{\max_{j\notin S}|Z_j|\leq t\mid X_S,w\}
 \leq\exp[-2(p-k)\overline\Phi(t)]\longrightarrow0,
\]
since the exponent's magnitude is at least
$c p^{\delta/4}/\sqrt{\log p}$. The exactness condition fails.
\end{proof}

\begin{proof}[Proof of \cref{regression:easy-thresholds}(b)]
Use the cap $q=(1-\zeta)m_0$, $\zeta=1/\log p$, in
\eqref{regression:saturation}, and write $t=q/\|r_S\|$.
The fit event and the hypothesis give
$t^2\geq(2+\delta/2)\Lambda_\lambda$ eventually.
Because $\Lambda_\lambda\geq\frac12\log p$, define
\[
 \Pi=2p\phi(t)/t
 \leq C\frac{n}{\lambda t}(p\lambda/n)^{-\delta/4}
 \leq C\frac n\lambda p^{-\delta/8}
 \leq C\sqrt n\,p^{-\delta/8}.
\]
Conditional on $(X_S,w)$, \eqref{regression:truncated-tail}
and the binomial bound imply
\begin{equation}\label{regression:cap-count}
 \sum_{j\notin S}(|Z_j|-t)_+^2\leq4(\Pi+\log p),\qquad
 |V|\leq2\Pi+3\log p=o(\sqrt n).
\end{equation}
Also $\max_{j\notin S}|Z_j|\leq2\sqrt{\log p}$ with probability
tending to one.

Decompose each null column as
$x_j=(a_j/\|r_S\|^2)r_S+x_j^\perp$.
Conditional on $(X_S,w,(a_j)_{j\notin S})$, the perpendicular
components are independent centered Gaussian vectors on $r_S^\perp$.
Thus choosing $V$ by its correlations does not bias these components.
Let $P$ project onto
$(\operatorname{range}X_S+\operatorname{span}r_S)^\perp$.
This space has dimension $n-k-1$ almost surely. On it,
$M_S^{-1}$ is the identity; consequently
$H=X_V^\top M_S^{-1}X_V\succeq X_V^\top PX_V$.
Conditional singular-value concentration and
\eqref{regression:cap-count} give
\[
 \lambda_{\min}(H)\geq n(1-o(1)),\qquad
 \|X_V^\perp\|_{\mathrm{op}}\leq\sqrt n(1+o(1)).
\]
The parallel component has squared operator norm at most
$4|V|\log p=o(n)$, so $\|X_V\|_{\mathrm{op}}\leq\sqrt n(1+o(1))$.
The construction is therefore defined and
\begin{equation}\label{regression:gamma-est}
 \Gamma\leq\frac{4\|r_S\|^2(\Pi+\log p)}{n(1-o(1))},\qquad
 \gamma_1:=\frac{\sqrt\Gamma}{\|r_S\|}
          =o(1/\log p).
\end{equation}

For $j\in V$, the new absolute correlation equals $q$.
For $j\notin S\cup V$, condition also on the selected perpendicular
columns, which determine $\Delta$. The remaining perpendicular
columns are still independent Gaussian vectors. A union bound gives
$|(x_j^\perp)^\top\Delta|\leq2\sqrt{\Gamma\log p}$ simultaneously.
It follows that
\[
 M\leq q(1+\gamma_1)+2\sqrt{\Gamma\log p}
       \leq m_0(1-\zeta+3\gamma_1),
\]
using $m_0/\|r_S\|\geq t$ and $t^2\geq\log p$.

For the selected correlations,
\[
 X_S^\top\Delta
 =\lambda(G+\lambda I)^{-1}X_S^\top X_Vu.
\]
The spectral estimates above and
$\|u\|\leq2\|r_S\|\sqrt{\Pi+\log p}/[n(1-o(1))]$ imply
\[
 \frac{\max_{i\in S}|x_i^\top\Delta|}{m_0}
 \leq C\left[
 \frac{\sigma}{b}\sqrt{\frac{\Pi+\log p}{n}}+
 \sqrt{\frac kn}\frac{\lambda}{n}\sqrt{\Pi+\log p}\right]
 =o(1/\log p).
\]
For the last equality, the $\log p$ terms use
$n\geq\log^6p$, $k\log p/n\leq C_0$, and
$\lambda/n\leq\log^{-2}p$. For the $\Pi$ terms use
$\Pi\leq C(n/\lambda)p^{-\delta/8}$ and $\lambda\geq\sqrt n$.
Therefore
\[
 m^2-M^2\geq(2-o(1))\zeta m_0^2.
\]
By \eqref{regression:gamma-est},
\[
 \frac{\Gamma}{(m^2-M^2)/\lambda}
 \leq(1+o(1))\frac{4\lambda(\Pi+\log p)}
                         {n\zeta t^2}\longrightarrow0.
\]
Indeed $\zeta t^2$ is bounded away from zero,
$\lambda\Pi/n=O(p^{-\delta/8})$, and
$\lambda\log p/n\leq1/\log p$.
\cref{regression:capped} now gives every wrong-fixing bound at least
$F_\lambda(S)+(2-o(1))\lambda b^2/\log p$.
This proves the stated margin and uniqueness.
\end{proof}

We first establish a deterministic feasible point for the negative
single-variable statement.

\begin{lemma}[A forced-in primal point]\label{regression:forced-primal}
Fix $S$ of size $k$ and $j\notin S$. Put
$b_+=\max_{i\in S}|\widehat\beta_i|$ and
$a=X^\top r_S$. Choose $V'\subseteq[p]\setminus(S\cup\{j\})$,
$q\geq\lambda b_+$, and define
\[
 e_l=(|a_l|-q)_+,\quad Q_e=\sum_{l\in V'}e_l^2,\quad
 A=\|X_{V'}\|_{\mathrm{op}}^2,\quad
 v_l=s\,\operatorname{sign}(a_l)e_l/A,\quad
 t_l=\lambda|v_l|/q,\quad T=\sum_lt_l .
\]
If $0<s\leq1$, $t_l\leq1$, and $T<k-1$, then
\begin{equation}\label{regression:forced-primal-eq}
 r(\varnothing,\{j\})\leq F_\lambda(S)+
 \lambda b_+^2\left[1+\frac{(1+T)^2}{k-1-T}\right]
                         -(2s-s^2)Q_e/A.
\end{equation}
\end{lemma}
\begin{proof}
Take $z_i=1-(1+T)/k$ on $S$, $z_j=1$, $z_l=t_l$ on $V'$,
and zero otherwise. Keep $\widehat\beta$ on $S$, take coefficients
$v$ on $V'$, and zero at $j$. The budget is exactly $k$.
Substitution in \eqref{regression:primal} bounds the objective by
\[
 F_\lambda(S)+\lambda\|\widehat\beta\|^2\frac{1+T}{k-1-T}
 -2\sum_l|v_l||a_l|+\|X_{V'}v\|^2+q\sum_l|v_l|.
\]
Use $\|\widehat\beta\|^2\leq kb_+^2$ and
$k(1+T)/(k-1-T)=1+T+(1+T)^2/(k-1-T)$.
The term $\lambda b_+^2T$ is at most $q\sum_l|v_l|$.
The last terms are consequently at most
$-2\sum_l|v_l|e_l+A\|v\|^2=-(2s-s^2)Q_e/A$.
\end{proof}

\begin{proof}[Proof of \cref{regression:easy-thresholds}(c)]
Let $d_0=\delta/16$, $q=\lambda b_+(1+1/\log p)$,
$t'=q/\|r_S\|$, and $t_0=(1+d_0)t'$.
On the fit event, $t_0^2\leq(2-\delta/2)\Lambda_\lambda$.
The number of nulls with $|a_l|/\|r_S\|\geq t_0$ has conditional
binomial mean $\mu$ satisfying
\[
 \mu\geq c\frac n\lambda
                 (p\lambda/n)^{\delta/4}/\sqrt{\log p}.
\]
Here one may use $1+t_0$ in the denominator of the elementary
Gaussian lower-tail bound; this also handles bounded $t_0$.
In particular $\mu/(n/\lambda)\to\infty$.
With probability tending to one, at least $\mu/2$ such nulls exist.
Let $N'=\min(\lfloor\mu/2\rfloor,\lfloor n/\log^2p\rfloor)$
and let $V'$ contain the $N'$ largest null absolute correlations.
This set is common to all forced-in indices $j\notin S\cup V'$.
On $V'$, $e_l\geq d_0q$.

Conditioning on all correlations leaves its perpendicular columns
independent Gaussian. Since $N'\leq n/\log^2p$ and
$\max_l|Z_l|\leq2\sqrt{\log p}$, the same decomposition as above gives
$A=\|X_{V'}\|_{\mathrm{op}}^2\leq n(1+o(1))$.
Both choices of $N'$ satisfy
$N'\geq3(1+o(1))n/(d_0^2\lambda)$ eventually:
the first because $\mu/(n/\lambda)\to\infty$, and the second because
$\sqrt n/\log^2p\to\infty$.
Thus $Q_e/A\geq3\lambda b_+^2$ for all sufficiently large $p$.
Set $s=3\lambda b_+^2A/Q_e\leq1$.
Then $(2s-s^2)Q_e/A\geq3\lambda b_+^2$ and
\[
 T=\frac{3\lambda^2b_+^2}{qQ_e}\sum_le_l\leq3/d_0.
\]
Also $\max_lt_l=o(1)$. To see this, $Q_e\geq N'd_0^2q^2$,
$\max e_l\leq2\|r_S\|\sqrt{\log p}$, and
$N'\gg\sqrt{\log p}$. The ratio $q/\|r_S\|$ is bounded below:
writing the square of \eqref{regression:tau} as a harmonic mean,
\[
 \tau_\lambda^{-2}
 =\frac{\sigma^2(n+\lambda)^2}{\lambda^2b^2n}+\frac kn,
\]
shows a fixed positive lower bound for $\tau_\lambda$ under
\eqref{regression:easy-regime}. The fit estimate transfers it to
$q/\|r_S\|$.

Since $k\geq5000/\delta^2$,
$(1+T)^2/(k-1-T)\leq1/2$; indeed
$2(1+48/\delta)^2+1+48/\delta\leq4851/\delta^2$.
\cref{regression:forced-primal} gives the node bound at most
$F_\lambda(S)-\tfrac32\lambda b_+^2$.
The coefficient estimate implies $b_+=b(1+o(1))$, proving the
claimed strict inequality with margin $\lambda b^2/2$.
There are at most $N'\leq n/\log^2p$ exceptional indices.
\end{proof}

\subsection{Ridge optimization}
\label{app:regression:sample-sizes}

\begin{proof}[Proof of \cref{regression:window}]
The preceding harmonic-mean identity shows exactly that
$\tau_\lambda^2<n/k$. Also, uniformly for the allowed ridges,
writing $\lambda=(\sigma/b)\sqrt{Tn}$ gives
\begin{equation}\label{regression:ridge-scaling}
 \tau_\lambda^2=(1+O(\log^{-2}p))\frac{Tn}{n+kT}.
\end{equation}
Let $\Lambda=\log p$ for the root test, and
$\Lambda=\log(p/\sqrt n)$ for the C1 test.
For any fixed $e\in(0,1)$, if $n\geq(2+e)k\Lambda$, set
\[
 A=(2+e/2)(1+e/8)\Lambda,\qquad
 T=\frac A{1-Ak/n},\qquad \lambda=(\sigma/b)\sqrt{Tn}.
\]
Because $(2+e/2)(1+e/8)<2+e$, the denominator is bounded away
from zero and $T=\Theta_e(\log p)$.
The inequalities $\log^6p\leq n\leq p$ imply that this ridge is
between $\sqrt n$ and $n/\log^2p$ eventually.
In the C1 case,
$\Lambda_\lambda=\Lambda+\frac12\log(T\sigma^2/b^2)
=(1+o(1))\Lambda$. Formula \eqref{regression:ridge-scaling}
then gives either positive conclusion of
\cref{regression:easy-thresholds} with a smaller fixed slack
than $e$, for example $e/4$ eventually. It need not meet the
threshold hypothesis with the original slack $e$.

If $n\leq(2-e)k\Lambda$, the exact inequality
$\tau_\lambda^2<n/k$ gives the negative root comparison.
For C1, also use $\Lambda_\lambda\geq\log(p/\sqrt n)$.
Apply the respective negative theorem with slack $e/2$ and its
corresponding lower bound on $k$.
For $k=p^{\gamma+o(1)}$ and $n=a k\log p$,
$\log(p/\sqrt n)=(1-\gamma/2+o(1))\log p$.
In $2-\gamma<a<2$, the positive C1 theorem proves $S^*$ uniquely
optimal, while the negative root theorem gives $R<F_\lambda(S^*)$.
The root is therefore globally inexact, and the path lemma gives
the tree bound. The other regimes follow in the same way.
\end{proof}

\subsection{Hull calculus and deterministic completion}
\label{app:regression:hull}

Zero extension maps generators of $\mathcal H_T$ to generators of a
larger hull. Moreover, for $N\succeq0$,
\begin{equation}\label{regression:hull-recession}
 (z,b,B)\in\mathcal H_T\quad\Longrightarrow\quad
                          (z,b,B+N)\in\mathcal H_T .
\end{equation}
For a rank-one $N=vv^\top$, mix $(z,b,B)$ with the generator
$(\mathbf1_T,v/\sqrt t,vv^\top/t)$ with weight $t$ and let
$t\downarrow0$. The first moments converge to $(z,b)$ and the
second moment to $B+vv^\top$. Closedness and a decomposition of
$N$ into rank-one terms prove \eqref{regression:hull-recession}.
The singleton hull is
$0\leq z\leq1$, $B\geq0$, $b^2\leq zB$.
Consequently the global positive semidefinite constraint and the
singleton hulls give, for $E=B-\beta\beta^\top$,
$E\succeq0$ and
$E_{jj}\geq\beta_j^2(1/z_j-1)$. This proves
$\ell_s(z)\geq g_\lambda(z)$; integer generators prove validity.
Convexity follows by projecting a convex feasible epigraph.

For a block $T$ and a positive semidefinite block matrix $Q_T$,
every generator satisfies
$b^\top Q_Tb=\langle Q_T,bb^\top\rangle$.
Convex combinations and limits imply that the exact closed
indicator epigraph hull has value at most $\langle Q_T,B_{TT}\rangle$
on $\mathcal H_T$. The residual matrix $Q_0\succeq0$ has
$\langle Q_0,B\rangle\geq\beta^\top Q_0\beta$.
Summing proves the selected-block dominance stated in the text.

\begin{lemma}[A support distribution]\label{regression:mixture}
For $z\in K$, with any prescribed zero and one fixings, there is a
random feasible binary support with indicator $\zeta$ and mean $z$.
For $\beta_j=0$ where $z_j=0$, define
$\xi_j=\beta_j\zeta_j/z_j$ when $z_j>0$ and zero otherwise.
Then $\E\xi=\beta$ and $D=\operatorname{Cov}\xi$ satisfies
$\operatorname{tr}D=\pi(z,\beta)$.
On each coordinate subset, $(z,\beta,\beta\beta^\top+D)$ restricts
to a point of its local hull.
\end{lemma}
\begin{proof}
Every vertex of $K$ is binary. Indeed, two fractional coordinates
can be perturbed in opposite directions, while a lone fractional
coordinate can be perturbed if the budget is slack. If the budget
is tight and only one coordinate is fractional, its value would
equal the difference of two integers, a contradiction.
The same argument applies to a face with fixings.
Express $z$ as a convex combination of its binary vertices.
The coefficient expectation and variance follow coordinatewise:
$\operatorname{Var}\xi_j=\beta_j^2(1/z_j-1)$.
The second-moment points are actual convex combinations of generators.
\end{proof}

\subsubsection{Completion using zero-indicator helper columns}
\label{app:regression:completion}

Let $F\supseteq\supp z$, and embed the covariance $D$ from
\cref{regression:mixture} in the $F$ block. Let $H$ be a disjoint
helper set and assume $W=X_HX_H^\top\succ0$. Define
\[
 P_H=I-X_H^\top W^{-1}X_H,\quad
 q_m=x_m^\top W^{-1}x_m,\quad
 \theta_F=\|X_F^\top W^{-1}X_F\|_{\mathrm{op}},\quad
 Y_F=X_FD^{1/2}.
\]
The matrix $P_H$ is the orthogonal projector onto $\ker X_H$.
For a slack $a>0$, set
\[
 G_F=\sqrt{1+a}\,D^{1/2},\qquad
 C=-X_H^\top W^{-1}X_FG_F,
\]
and form $G$ by placing $G_F$ on $F$, $C$ on $H$, and zeros
elsewhere. Then $XG=0$.

\begin{lemma}[Local completion]\label{regression:completion}
Let $J\succeq0$ have range in $\ker X_H$ and be embedded in the
helper block. Put $B=\beta\beta^\top+GG^\top+J$.
Its lifted objective is
\begin{equation}\label{regression:completion-cost}
 \mathcal F(\beta,B)=
 \|y-X\beta\|^2+\lambda\|\beta\|^2+
 \lambda[(1+a)\pi+\|C\|_F^2+\operatorname{tr}J],
 \quad
 \|C\|_F^2\leq(1+a)\theta_F\pi.
\end{equation}
Two choices ensure local hull feasibility:
\begin{enumerate}[label=(\alph*)]
\item For hulls of order $s$, suppose
 $\eta_H=\max_{M\subseteq H,\ |M|\leq s-1}
             \|X_M^\top W^{-1}X_M\|_{\mathrm{op}}<1$.
 Set $J=tP_H$, where
 \[
 t=\frac{1+a}{a(1-\eta_H)}
           \max_{|M|\leq s-1}\|C_M\|_{\mathrm{op}}^2 .
 \]
 Then $B$ is feasible for all hulls of order at most $s$, and
 $\operatorname{tr}J=t(|H|-n)$. For $s=1$ take $J=0$.
\item For pairwise hulls, suppose explicitly that
 $q_{\max}=\max_{m\in H}q_m<1$.
 Set
 \[
 J=P_H\operatorname{diag}
       \left(\frac{\|C_m\|^2}{a(1-q_m)^2}:m\in H\right)P_H.
 \]
 Then $B$ is feasible for pairwise hulls and
 \[
 \operatorname{tr}J\leq\frac{\|C\|_F^2}{a(1-q_{\max})}.
 \]
 Thus the additional cost above the perspective objective is at most
 $\lambda\pi[a+(1+a)\theta_F(1+1/[a(1-q_{\max})])]$.
\end{enumerate}
\end{lemma}
\begin{proof}
Positive semidefiniteness is immediate from $GG^\top+J\succeq0$.
Since $XG=0$ and $X_HJ=0$, its design-variance cost vanishes.
The trace is $(1+a)\operatorname{tr}D+\|C\|_F^2+
\operatorname{tr}J$, giving \eqref{regression:completion-cost}.
Also
$\|C\|_F^2=(1+a)\operatorname{tr}(Y_F^\top W^{-1}Y_F)
\leq(1+a)\theta_F\pi$.

For a tested subset $T$, put $T'=T\cap F$ and $M=T\cap H$.
The random-support moment on $T'$ belongs to $\mathcal H_{T'}$.
By zero extension and \eqref{regression:hull-recession}, it is
enough that the additional block
\[
 \begin{pmatrix}
 C_MC_M^\top+J_{MM}&C_MG_F^\top\\
 G_FC_M^\top&aD
 \end{pmatrix}
\]
is positive semidefinite, before restriction to $T'$.
If either block is absent, this is immediate.
For (a), $J_{MM}\succeq t(1-\eta_H)I$.
If $C\ne0$ and $t>0$, its Schur complement is bounded by
$\|C_M\|^2(1+a)D/[t(1-\eta_H)]\preceq aD$.
If the maximum in $t$ is zero, every relevant $C_M$ vanishes;
the blocks are then diagonal and positive semidefinite.
This handles the zero case without invoking a positive definite
helper covariance.

In (b), a mixed tested pair has $M=\{m\}$ and
$J_{mm}\geq\|C_m\|^2/a$, because the diagonal contribution of
$P_H$ is $1-q_m$. The Schur complement is at most
$(1+a)\|C_m\|^2D/(\|C_m\|^2+J_{mm})\preceq aD$.
If $C_m=0$, the cross block vanishes.
Finally, $P_H^2=P_H$ gives
$\operatorname{tr}J=\sum_m\|C_m\|^2/[a(1-q_m)]$.
\end{proof}

\subsection{Proof of the local-lift easy thresholds}
\label{app:regression:lift-easy}

Split the null features, in a deterministic order, into two halves
$H_1,H_2$. For a helper half $H$, all constructions below are
measurable with respect to $\mathcal A_H=\sigma(X_{[p]\setminus H},w)$.
The helper matrix is independent of this information.
Let $m=|H|$, and define
\[
 w_-=(\sqrt m-\sqrt n-2\sqrt{\log p})^2,\quad
 v_+=n+2\sqrt{3n\log p}+6\log p,\quad
 a_p=1+2\sqrt{3\log p}+6\log p .
\]
Under $s n\log p=o(p)$, the right side before squaring in $w_-$
is positive eventually, and $w_-=m(1-o(1))$.
With probability tending to one,
$\lambda_{\min}(W)\geq w_-$ and $\max_{j\in H}\|x_j\|^2\leq v_+$.
For any fixed finite number of $\mathcal A_H$-measurable matrices $Y$,
simultaneously for every $j\in H$,
\begin{equation}\label{regression:helper-row}
 \|x_j^\top W_{-j}^{-1}Y\|^2
       \leq a_p\|W_{-j}^{-1}Y\|_F^2,\qquad
 W_{-j}=W-x_jx_j^\top .
\end{equation}
To prove this, condition on $(X_{H\setminus\{j\}},\mathcal A_H)$.
The remaining $x_j$ is standard Gaussian, so the weighted
chi-square bound, obtained by applying its exponential moment,
gives
$x_j^\top A x_j\leq\operatorname{tr}A+
2\sqrt{3\log p}\|A\|_F+6\log p\|A\|_{\mathrm{op}}$
for $A=W_{-j}^{-1}YY^\top W_{-j}^{-1}$.
The last expression is at most $a_p\operatorname{tr}A$.
A union bound proves \eqref{regression:helper-row}.

Sherman--Morrison, $W_{-j}\succeq(w_--v_+)I$, and
\eqref{regression:helper-row} imply
\[
 \max_{|M|\leq s-1}\|C_M\|^2
 \leq\frac{(1+a)(s-1)a_p\|Y_F\|_F^2}{(w_--v_+)^2},
 \qquad \eta_H\leq(s-1)v_+/w_- .
\]
If $\|X_F\|_{\mathrm{op}}^2\leq16n$, then
$\|Y_F\|_F^2\leq16n\pi$.
Choose the completion slack $a=\sqrt{s n\log p/p}$.
\cref{regression:completion}(a) gives
\begin{equation}\label{regression:inflation}
 \ell_s(z)\leq P_e(z,\beta):=
 \|y-X\beta\|^2+\lambda\|\beta\|^2+\lambda(1+e)\pi(z,\beta),
 \qquad e=O(\sqrt{s n\log p/p})=o(1).
\end{equation}
The constant is uniform for the finite number of matrices used
below. It does not require independence of $W$ and its own
individual columns: \eqref{regression:helper-row} used leave-one-out
conditioning.

\begin{proof}[Proof of \cref{regression:lift-thresholds}]
The positive assertions follow from perspective domination and validity.
For the negative root assertion, use $H_1$ as helpers and $H_2$ as
candidates. On the fit event a candidate $l\in H_2$ satisfies
$|a_l|\geq(1+c)m_0$ with probability tending to one, where
$c=\delta/12$ and $m_0=\lambda\min_S|\widehat\beta_i|$.
This follows from the Gaussian maximum lower bound at
$\sqrt{(2-\delta/2)\log p}$, whose ratio to the selected threshold
is eventually greater than $1+c$.
Choose $i\in S$ attaining $m_0$, set $\bar e=c/4$, and take
\[
 t=\min(c/8,c\lambda/(8v_+)),\quad
 v=t a_l/[\lambda(1+\bar e)],\quad
 z=\mathbf1_S-te_i+te_l,\quad \beta=\widehat\beta+ve_l.
\]
Use the two-support mixture on $S$ and $S\setminus\{i\}\cup\{l\}$.
All these objects are $\mathcal A_{H_1}$-measurable.
The global column-norm bound and the fixed $X_S$ spectral bound
give $\|X_{S\cup\{l\}}\|^2\leq16n$.
Eventually $e\leq\bar e$ in \eqref{regression:inflation}.
Direct substitution gives
\[
 P_{\bar e}(z,\beta)
 \leq F_\lambda(S)-\frac t\lambda
 \left[\frac{a_l^2}{1+\bar e}
       -\frac{(1+\bar e)m_0^2}{1-t}
       -\frac{t a_l^2\|x_l\|^2}{\lambda}\right].
\]
Dividing the bracket by $a_l^2$ bounds it below by
\[
 (1-c/4)
 -\frac{1+c/4}{(1+c)^2(1-c/8)}-c/8\geq c/8>0.
\]
Here $(1+c/4)/(1-c/8)\leq1+c/2$ and
$(1+c/2)/(1+c)^2\leq1/(1+c)\leq1-c/2$.
Thus $\mathcal R(\varnothing,\varnothing)<F_\lambda(S)$.

For the forced-in converse, repeat the construction of
\cref{regression:forced-primal} separately in each candidate half,
using the other half as helpers. Let $d_0=\delta/16$,
$q=\lambda b_+(1+1/\log p)$ and select the $N'$ largest
absolute correlations in the candidate half, with
$N'=\min(\lfloor\mu/2\rfloor,\lfloor n/\log^2p\rfloor)$
as before. Half the number of candidates changes no limiting estimate.
For any nonexceptional $j$ in this half, put
$\bar e=d_0/2$ and $e_l'=(|a_l|-(1+\bar e)q)_+$ on $V'$.
Then $e_l'\geq(d_0/2)q$.
Keep the definitions of $v_l,t_l,T$ in the primal lemma, using
$e_l'$ in place of $e_l$. Substitution in $P_{\bar e}$ gives
\begin{equation}\label{regression:inflated-forced}
 P_{\bar e}(z,\beta)\leq F_\lambda(S)
 +(1+\bar e)\lambda b_+^2
       \left[1+\frac{(1+T)^2}{k-1-T}\right]
                  -(2s_0-s_0^2)Q'/A,
\end{equation}
where $Q'=\sum_l(e_l')^2$ and
$v_l=s_0\operatorname{sign}(a_l)e_l'/A$.
For clarity, the ridge increment on $S$ is
$(1+\bar e)\lambda\|\widehat\beta\|^2(1+T)/(k-1-T)$,
and on $V'$ it is at most $(1+\bar e)q\sum_l|v_l|$.
The same rearrangement as in \cref{regression:forced-primal}
then gives \eqref{regression:inflated-forced}.

Use a support mixture that includes $j$ always and chooses a
support of size at most $k-1$ on $S\cup V'$. Its coefficient on
$j$ is zero, so its covariance has a zero $j$ row and column.
Set $F=S\cup V'\cup\{j\}$, padding $D$ with that zero row.
The matrix $Y_F=X_FD^{1/2}$ is the same for all nonexceptional $j$
in this half, and is independent of the helper half.
The norm of $X_F$ is bounded uniformly by $16n$ in squared
operator norm: the bounds for $X_S$ and $X_{V'}$, together with
the global column-norm bound for $x_j$, suffice.
Thus \eqref{regression:inflation} applies simultaneously and
$\mathcal R(\varnothing,\{j\})\leq P_{\bar e}(z,\beta)$.

As in the perspective proof, eventually
$N'\geq12(1+o(1))n/(d_0^2\lambda)$.
Take $s_0=3\lambda b_+^2A/Q'\leq1$.
Then $T\leq6/d_0=96/\delta$, and $\max_l t_l=o(1)$.
The bound $k\geq20000/\delta^2$ ensures
$(1+T)^2/(k-1-T)\leq1/2$. Formula
\eqref{regression:inflated-forced} is consequently at most
$F_\lambda(S)+\frac32(1+\bar e)\lambda b_+^2
-3\lambda b_+^2<F_\lambda(S)-\lambda b_+^2$.
This proves the desired margin, with at most $2N'$ exceptions
over the two halves. The sample-size statements follow from
\cref{app:regression:sample-sizes}, again with a smaller fixed
threshold slack.
\end{proof}

\subsection{Conflict packing and the retained-helper hard theorem}
\label{app:regression:hard}

\begin{lemma}[Integer optimum lower bound]\label{regression:opt-lower}
Let $q=k/n<1$, and let $\rho\in(q,1)$ satisfy
\[
 \frac n2\left[q\log(q/\rho)+(1-q)\log((1-q)/(1-\rho))\right]
       =\log\binom pk+\log(1/\alpha).
\]
If $y$ is independent of $X$, then with probability at least
$1-\alpha$, $\OPT\geq(1-\rho)\|y\|^2$.
In the planted model the same probability bound holds for
$\OPT\geq(1-\rho)\sigma^2\|P_{S^*}^\perp w\|^2$.
\end{lemma}
\begin{proof}
For a fixed $k$-subset $T$, the span of its columns is a uniform
$k$-dimensional subspace. Conditional on an independent nonzero
vector $v$, the squared projection fraction has the law
$A/(A+C)$, where $A\sim\chi_k^2$ and $C\sim\chi_{n-k}^2$ are
independent. For $\rho>q$, exponential Markov applied to
$(1-\rho)A-\rho C$ gives
\[
 \Pbb\{A/(A+C)\geq\rho\}\leq
 (1-2t(1-\rho))^{-k/2}(1+2t\rho)^{-(n-k)/2}.
\]
Choosing $2t(1-\rho)=1-q/\rho$ yields the exponential of minus
the left side of the displayed equation. Union-bound over all
$k$-subsets. Ridge values dominate projection residuals, and
padding any smaller support to size $k$ cannot increase its
value. This proves the independent case.

For the planted case, for each $T$ let
$v_T=X_{S^*\setminus T}\beta_{S^*\setminus T}^*+\sigma w$.
It is independent of $X_T$, and $y-v_T$ lies in its span.
Therefore the same fixed-support probability bound applies.
Moreover $\|v_T\|\geq\sigma\|P_{S^*}^\perp w\|$.
The union bound again gives the result. Independence of the
different residuals is unnecessary.
\end{proof}

\begin{proof}[Proof of \cref{regression:hard-theorem}]
Put $\alpha=1/k$. Since $k/n\to0$ and $x_p\to x$,
the solution $\rho$ in \cref{regression:opt-lower} satisfies
$1-\rho=e^{-x}+o(1)$. This follows from
$q\log(q/\rho)+(1-q)\log((1-q)/(1-\rho))
=-\log(1-\rho)+o(1)$, uniformly on compact subsets of $(0,1)$.
In case (a),
\begin{equation}\label{regression:hard-opt}
 \OPT\geq(e^{-x}-o(1))\|y\|^2
\end{equation}
with probability tending to one.

Choose constants $\mu,\vartheta\in(0,1)$ and
$0<c<(1-\mu)\vartheta$ such that, with
$B_0=1+2\sqrt{cx}+2cx$,
\begin{equation}\label{regression:hard-parameters}
 \frac{(1+\mu)(1-\vartheta)x}{B_0}>e^x-1 .
\end{equation}
These exist because $2x>e^x-1$: let $\mu$ tend to one and
$\vartheta,c$ tend to zero, then use a strict margin.
Write $M=\lceil k(p/k)^\vartheta\rceil$ and
$m=\lceil\mu k\rceil$. Conditional on $y$, the correlations
$c_j=x_j^\top(y/\|y\|)$ are independent standard normal variables.
The number whose absolute values exceed
$t_M=\overline\Phi^{-1}(M/p)$ has mean $2M$, and is at least
$M$ with probability at least $1-e^{-M/4}$.
Let $W_0$ be the $M$ largest absolute correlations.
The binomial coefficient estimates
$(p/k)^k\leq\binom pk\leq(ep/k)^k$ imply
\[
 \log(p/k)=\frac{x+o(1)}2\,\frac nk\longrightarrow\infty,
 \qquad
 t_M^2=2(1-\vartheta)\log(p/k)(1+o(1)).
\]
Every $U\subseteq W_0$ with $|U|\geq k+m$ therefore has
\begin{equation}\label{regression:hard-correlation}
 s_U:=\sum_{j\in U}c_j^2
       \geq(1+\mu)(1-\vartheta)xn(1+o(1)).
\end{equation}

Greedily pack the $k$-subsets of $W_0$ in a fixed lexicographic
order so that $|S\setminus T|\geq m$ for distinct members.
Each chosen support deletes at most
$\sum_{d<m}\binom kd\binom{M-k}d$ candidates, so the family has
size at least the total number of candidates divided by this sum.
Using $\binom kd\leq2^k$ and
$\binom{M-k}d\leq(eM/d)^d$,
\[
 \log|\mathcal C_0|
 \geq(1-\mu)\vartheta k\log(p/k)-O(k)-O(\log(p/k)).
\]
Since $k\to\infty$, select the first
$\lceil\exp(c k\log(p/k))\rceil$ members to obtain $\mathcal C$.
This code is measurable in the correlations before the
perpendicular components of the columns are revealed.

Conditional on $y$ and all $c_j$,
$X^\perp=(I-\widehat y\widehat y^\top)X$ has independent
Gaussian columns on $\widehat y^\perp$, where
$\widehat y=y/\|y\|$. For a fixed union $U$,
$\|X_U^\perp c_U\|^2/s_U$ has law $\chi_{n-1}^2$.
Union-bound over the at most $|\mathcal C|^2$ pairs, with
$t=\log(|\mathcal C|^2/\alpha)=cxn(1+o(1))$ in
\eqref{regression:chi-tail}. With probability tending to one,
all pairs satisfy
\begin{equation}\label{regression:hard-perp}
 \|X_U^\perp c_U\|^2/s_U\leq nB_0(1+o(1)).
\end{equation}
No independence between different unions is used.

For $z=(\mathbf1_S+\mathbf1_T)/2$, perspective fitting on
$U=S\cup T$ has ridge coefficient $\lambda$ on the intersection
and $2\lambda$ on the symmetric difference. It is bounded above
by the fit with coefficient $2\lambda$ on all of $U$.
Taking $\beta=t c_U$ and minimizing this scalar $t$ gives
\begin{equation}\label{regression:hard-mid}
 g_\lambda(z)\leq
 \|y\|^2\left[1-\frac{s_U}
 {s_U+\|X_U^\perp c_U\|^2/s_U+2\lambda}\right].
\end{equation}
The identity follows by writing
$X_U=\widehat y c_U^\top+X_U^\perp$ and expanding the squared loss.
Since $|U|\geq k+m$ and $\lambda=o(n)$,
\eqref{regression:hard-correlation}--\eqref{regression:hard-mid}
give, uniformly over pairs,
\[
 g_\lambda(z)\leq\|y\|^2
 \left[\frac{B_0}{B_0+(1+\mu)(1-\vartheta)x}+o(1)\right].
\]
The coefficient in brackets is strictly below $e^{-x}$ by
\eqref{regression:hard-parameters}. Together with
\eqref{regression:hard-opt}, this leaves a fixed positive
margin. Choose $\eta$ strictly smaller than that margin.

For case (b), perform the correlation and code construction on
$[p]\setminus S^*$, conditionally on $(X_{S^*},w)$.
There are $p-k=p(1-o(1))$ candidates, so all estimates are unchanged.
The optimum bound is now
$\OPT\geq(e^{-x}-o(1))n\sigma^2$, while
$\|y\|^2=(1+\kappa+o_{\Pbb}(1))n\sigma^2$.
Choose the parameters instead to satisfy
\[
 (1+\kappa)\frac{B_0}{B_0+(1+\mu)(1-\vartheta)x}<e^{-x}.
\]
They exist because $\kappa<e^{-x}(1+2x)-1$.
This proves the perspective conflicts in both cases.

It remains to justify the same conflicts for $\ell_2$ uniformly.
The hard hypotheses imply $\log p=o(n)$ and $p/n\to\infty$.
Indeed $\log(p/k)\asymp n/k$, $\log k=o(n)$, and
$p/n=(p/k)/(n/k)\to\infty$.
There is a simultaneous event of probability tending to one on which
\begin{equation}\label{regression:uniform-hard-design}
 \lambda_{\min}(XX^\top)=p(1-o(1)),\quad
 \max_j\|x_j\|^2=n(1+o(1)),\quad
 \max_{|U|\leq2k}\|X_U\|_{\mathrm{op}}^2=O(n).
\end{equation}
The first two follow from
\eqref{regression:singular-tail} applied to $X^\top$ and from
\eqref{regression:chi-tail}, with union bounds.
For the third, pad to $h=2k$ columns and apply the largest
singular-value bound with deviation
$\sqrt{4h\log(ep/h)}$. Union-bound over at most
$\exp(h\log(ep/h))$ subsets. The failure is at most
$\exp(-h\log(ep/h))$, and the squared bound
$(\sqrt n+\sqrt h+\sqrt{4h\log(ep/h)})^2$ is $O(n)$.

For every such $U$, use $H=[p]\setminus U$ as helpers.
On \eqref{regression:uniform-hard-design},
\[
 \lambda_{\min}(X_HX_H^\top)
 \geq\lambda_{\min}(XX^\top)-\|X_U\|^2=p(1-o(1)).
\]
Consequently $\theta_U=O(n/p)$ and $q_{\max}=O(n/p)<1$,
uniformly in $U$. Apply \cref{regression:completion}(b) with
$a=\sqrt{n/p}$.
For the two-support mixture with mean
$z=(\mathbf1_S+\mathbf1_T)/2$ and any coefficient supported on
$U$, $\pi(z,\beta)\leq\|\beta\|^2$. The completion therefore gives
\[
 \ell_2(z)\leq\|y-X_U\beta\|^2+(2+e_p)\lambda\|\beta\|^2,
 \qquad e_p=O(\sqrt{n/p})=o(1),
\]
simultaneously for every code pair. All support-selection dependence
is handled by the deterministic uniform event; the helpers need
not be independent of the selected union.
Taking again $\beta=t c_U$ replaces $2\lambda$ in
\eqref{regression:hard-mid} by $(2+e_p)\lambda=o(n)$.
The same strict margin proves the pairwise-hull conflicts.
Shrinking $\eta$ once gives the common assertion.

Convexity of a leaf containing two code points forces their midpoint
into that leaf, contradicting its certificate bound. Charging
certified removals as additional convex pieces gives at least
$|\mathcal C|$ terminal or removed pieces. For binaries each node
can fix each coordinate at most once, hence creates at most $p$
such removed pieces; leaves plus removals are at most $(p+1)$
times the node count. The lifted construction retains its helper
second moments at zero fixings. Deleting those columns invalidates
this feasible point and is not a consequence of the argument.
\end{proof}

\subsection{Variance price for the richer moment lift}
\label{app:regression:variance}

\begin{proof}[Proof of \cref{regression:variance-price}]
Represent the positive semidefinite matrix $Y$ as the Gram
matrix of $v_0,v_{z_j},v_{\beta_j}$, with $\|v_0\|=1$.
Set $w_j=v_{z_j}-z_jv_0$ and
$u_j=v_{\beta_j}-\beta_jv_0$.
Then $B-\beta\beta^\top=\operatorname{Gram}(u)$.
Let $F=\{j:0<z_j<1\}$ and let $V$ be the span of its $w_j$.
For $m\notin A=\supp z$, positive semidefiniteness and
$U_{mm}=\beta_m$ give $\beta_m^2\leq z_mB_{mm}=0$.
The product cones and $Z_{jm}\leq z_m=0$ give $U_{jm}=0$
for $j\ne m$. Therefore
$\langle u_m,w_j\rangle=U_{jm}-z_j\beta_m=0$:
the zero-indicator coefficient vectors are orthogonal to $V$.

For $j\in F$,
$\langle u_j,w_j\rangle=\beta_j(1-z_j)$ and
$\|w_j\|^2=z_j(1-z_j)$.
Put $g_j=P_Vu_j$ for every coordinate.
Orthogonal decomposition gives
$\operatorname{Gram}(u)\succeq\operatorname{Gram}(g)$,
and $g_m=0$ off $A$.
Thus
\[
 \langle Q,B-\beta\beta^\top\rangle
 \geq\langle X_A^\top X_A+\lambda I,\operatorname{Gram}_A(g)\rangle
 \geq(\nu_A+\lambda)\sum_{j\in F}\|g_j\|^2.
\]
Cauchy--Schwarz yields
$\|g_j\|^2\geq\beta_j^2(1-z_j)/z_j$.
Summing and adding the deterministic coefficient-fit objective
proves \eqref{regression:variance-price-eq}. Coordinates with
$z_j=1$ contribute zero to $\pi$, so none are missing.
\end{proof}

\subsection{The fixed-dimensional per-entry-noise limit}
\label{app:regression:pwe}

\begin{proof}[Proof of \cref{regression:pwe-theorem}]
Here $d,k,\beta^*,\sigma$ are fixed and $\lambda=\sqrt n$.
Finite-dimensional laws of large numbers give
$X^\top X/n\to I_d$, $X^\top w/n\to0$, and $\|w\|^2/n\to1$
in probability. For every $T$ with $|T|\leq k$,
\[
 \frac{F_\lambda(T)}n=
 \frac{\|y\|^2}n-
 \left(\frac{X_T^\top y}n\right)^\top
 \left(\frac{X_T^\top X_T}n+\frac{\lambda}nI\right)^{-1}
 \left(\frac{X_T^\top y}n\right)
 \longrightarrow \sigma^2+\|\beta^*_{S\setminus T}\|^2.
\]
For the empty support the subtraction is zero.
There are finitely many supports; convergence holds simultaneously.
Every $T\ne S$ in this family misses at least one selected coordinate
and has limiting gap at least $b_{\min}^2$ from $S$.
Let $U_n$ be the event that $S$ is uniquely optimal.
Then $\Pbb(U_n)\to1$. Each restricted ridge fit is unique, and on
$U_n$ its coefficients on $S$ are nonzero, since a zero coefficient
would give the same value on a proper subset. Thus the global
cardinality-constrained coefficient vector is unique as well.

Write $C_n=X_S^\top X_S$, $A_n=C_n+\sqrt nI$, and
$r=M_S^{-1}y$. Conditional on $X_S$,
\[
 \widehat\beta_S\sim
 N(A_n^{-1}C_n\beta_S^*,
       \sigma^2A_n^{-1}C_nA_n^{-1}).
\]
The mean tends to $\beta_S^*$ and $n$ times the covariance tends
to $\sigma^2I_k$. In particular $\widehat\beta_S\to\beta_S^*$.
The exact normal-equation identity gives
\[
 m_n:=\min_{i\in S}|x_i^\top r|
       =\sqrt n\min_{i\in S}|\widehat\beta_i|,
 \qquad m_n/\sqrt n\longrightarrow b_{\min}.
\]
For the residual, its signal part is
$q_n=\sqrt nX_SA_n^{-1}\beta_S^*$ and
\[
 \|q_n\|^2=n(\beta_S^*)^\top A_n^{-1}C_nA_n^{-1}\beta_S^*
                       \longrightarrow\|\beta_S^*\|^2.
\]
The matrix $M_S^{-1}-I$ acts only on the fixed $k$-dimensional
column space and has norm at most one.
Conditional on $X_S$, $\|P_Sw\|^2\sim\chi_k^2$.
Consequently $\|r-\sigma w\|=O_{\Pbb}(1)$, and
\[
 \frac{\|r\|^2}n\longrightarrow\sigma^2,\qquad
 \frac{m_n}{\|r\|}\longrightarrow b_{\min}/\sigma
\]
in probability.

Compute the planted-support exactness probability before conditioning
on its optimality. Conditional on $\mathcal A_n=\sigma(X_S,w)$,
the null columns are independent standard Gaussian vectors, whereas
$r$ and $m_n$ are fixed. The residual is nonzero almost surely,
since $M_S$ is invertible and the noise has a density. Therefore
the exact condition \eqref{regression:exact-condition-eq} gives
\[
 \Pbb\{R=F_\lambda(S)\mid\mathcal A_n\}
   =\left[1-2\overline\Phi(m_n/\|r\|)\right]^{d-k}.
\]
Equalities between a null absolute correlation and $m_n$ have
conditional probability zero. Thus the weak inequality used in
the exact convex first-order condition and the strict selected
inequality in the published certificate give the same probability.
The displayed random variable is bounded and converges in
probability to the right side of \eqref{regression:pwe-limit},
so its expectations converge to that constant.
This calculation uses conditional independence; the unscaled
null correlations share a random variance and need not be
unconditionally independent.

Finally, always $R\leq\OPT\leq F_\lambda(S)$.
The event $R=F_\lambda(S)$ implies $R=\OPT$, and on $U_n$
the two events coincide. Hence
\[
 0\leq\Pbb\{R=\OPT\}-\Pbb\{R=F_\lambda(S)\}
                             \leq\Pbb(U_n^c)\longrightarrow0.
\]
This proves the global exactness limit without conditioning the
Gaussian correlations on an event involving the null columns.
\end{proof}
